% Part: history
% Chapter: biographies
% Section: gerhard-gentzen

\documentclass[../../../include/open-logic-section]{subfiles}

\begin{document}

\olfileid{his}{bio}{gen}

\olsection{گرهارت گنتسن}

\olphoto{gentzen-gerhard}{Gerhard Gentzen}

گرهارت گنتسن بیش از همه به‌عنوان بنیان‌گذار نظریهٔ ساختاری اثبات، و
به‌ویژه به سبب ابداع دستگاه‌های !!{derivation} استنتاج طبیعی و حساب
دنباله‌ای، شناخته می‌شود. او در 24 نوامبر 1909 در گرایفسوالد آلمان به
دنیا آمد. گرهارت پیش از ورود به مدرسهٔ مقدماتی سه سال در خانه آموزش
دید و هنگام ورود از نظر تحصیلی از بیشتر هم‌کلاسی‌هایش عقب‌تر بود.
بااین‌حال دانش‌آمندی درخشان بود و استعداد چشمگیری در ریاضیات نشان داد.
علایق او گوناگون بود؛ برای نمونه، برای مادرش شعر و برای تئاتر مدرسه
نمایش‌نامه نیز می‌نوشت.

گنتسن تحصیلات دانشگاهی را در دانشگاه گرایفسوالد آغاز کرد، اما میان
گوتینگن، مونیخ و برلین جابه‌جا شد. او در سال 1933 زیر نظر هرمان وایل از
دانشگاه گوتینگن دکتری گرفت. پل برنایز بر بیشتر کار او نظارت داشت، اما
نازی‌ها برنایز را از دانشگاه اخراج کردند. گنتسن در سال 1934 کار خود را
به‌عنوان دستیار داوید هیلبرت آغاز کرد. همان سال دستگاه‌های
!!{derivation} حساب دنباله‌ای و استنتاج طبیعی را در مقاله‌های
\emph{Untersuchungen \"{u}ber das logische Schlie\ss en I--II
  [Investigations Into Logical Deduction I--II]} پدید آورد. او در سال
1936 سازگاری اصول موضوعهٔ پئانو را اثبات کرد.

رابطهٔ گنتسن با نازی‌ها پیچیده است. همان زمانی که برنایز، استاد
راهنمایش، ناچار به ترک آلمان شد، گنتسن به شاخهٔ دانشگاهی اس‌آ، سازمان
شبه‌نظامی نازی، پیوست. او مانند بسیاری از آلمانی‌ها عضو حزب نازی بود.
در دوران جنگ به‌عنوان افسر مخابرات در واحد اطلاعات هوایی خدمت کرد،
اما در سال 1942 به سبب فروپاشی عصبی از خدمت مرخص شد. روشن نیست وفاداری
گنتسن با حزب نازی بود یا برای تضمین موفقیت دانشگاهی به حزب پیوست.

در سال 1943 مؤسسهٔ ریاضیات دانشگاه آلمانی پراگ جایگاهی دانشگاهی به
گنتسن پیشنهاد کرد و او پذیرفت. اما در سال 1945 شهروندان پراگ علیه
اشغال آلمان شوریدند. نیروهای شوروی وارد شهر شدند و همهٔ استادان
دانشگاه را بازداشت کردند. گنتسن به سبب عضویت در سازمان‌های نازی به
اردوگاه کار اجباری فرستاده شد. او در 4 اوت 1945، در 35سالگی، بر اثر
سوءتغذیه در سلول خود درگذشت.

\begin{reading}
برای زندگی‌نامه‌ای کامل از گنتسن، \citet{Menzler-Trott2007} را ببینید.
شرحی خواندنی دربارهٔ ریاضی‌دانان زیر حاکمیت نازی که یادداشتی کوتاه نیز
دربارهٔ زندگی گنتسن دارد در \citet{Segal2014} آمده است. مقاله‌های
گنتسن دربارهٔ استنتاج منطقی در متن اصلی آلمانی در دسترس‌اند
\citep{Gentzen1935a,Gentzen1935b}. ترجمه‌های انگلیسی مقاله‌های گنتسن
را \citet{Gentzen1969} در یک مجلد گرد آورده است که طرحی زندگی‌نامه‌ای
را نیز در بر می‌گیرد.
\end{reading}

\end{document}
